\chapter{PENUTUP}
\label{sec:chap5_tutup}
\vspace{1ex}

Bab ini menyajikan kesimpulan yang ditarik dari hasil analisis pengujian sistem drone otonom pendeteksi manusia berbasis \textit{edge computing}, serta menjabarkan saran akademis dan teknis untuk pengembangan penelitian di masa mendatang.

\section{Kesimpulan}
\label{sec:sec5_kesimpulan}
\vspace{1ex}

Berdasarkan perancangan, implementasi, serta serangkaian pengujian terintegrasi yang telah dilakukan, dapat ditarik beberapa kesimpulan utama untuk menjawab rumusan masalah dan tujuan penelitian:
\begin{enumerate}
	\item Integrasi \textit{Edge Computing} Otonom: Sistem drone untuk misi Pencarian dan Penyelamatan (SAR) telah diintegrasikan pada \textit{companion computer} Raspberry Pi 5. Pendekatan optimasi—meliputi pemakaian sistem operasi \textit{headless}, konversi bobot model ke format ONNX, dan pemilihan resolusi $320\times320$ piksel—dapat meminimalkan beban komputasi. Konfigurasi ini mendukung kelancaran \textit{state machine} navigasi (transisi keadaan SCAN, APPROACH, DOCUMENT, dan DISPLACING) pada frekuensi kontrol 20 Hz dengan rata-rata latensi \textit{end-to-end} sebesar 188,04 ms di mana 95\% data berada di bawah 208,27 ms.
	\item Logika \textit{Persistent Multi-Target Scouting}: Sistem mampu menjalankan misi penjelajahan mandiri tanpa campur tangan operator melalui mode \textit{Scout Mission}. Pemanfaatan memori target historis (\textit{duplicate skip}) yang dipadukan dengan perhitungan radius eksklusi geospasial berbasis jarak \textit{Haversine} memungkinkan drone melakukan pemindaian, pendekatan, dokumentasi, dan perpindahan area secara berkelanjutan untuk mengidentifikasi beberapa target dalam satu sesi penerbangan.
	\item Sistem Keselamatan Terbang (\textit{Safety Procedure}): Arsitektur keselamatan tiga lapis (\textit{Three-Layer Safety System}) berfungsi sebagai jaring pengaman operasional drone. Validasi penerbangan nyata membuktikan bahwa manuver pendakian paksa (\textit{Pre-RTL Safety Climb}) dapat menaikkan ketinggian drone secara vertikal (dengan rentang kenaikan 0,46--2,27 meter) ketika kapasitas Battery RTL kritis terdeteksi, sebelum kendali autopilot PX4 mengambil alih secara penuh. Sistem juga menunjukkan keandalan dalam memotong loop navigasi otonom secara instan saat terjadi intervensi manual (\textit{RC Override}) atau ketika keluar dari mode OFFBOARD.
	\item Karakteristik Deteksi dan Keterbatasan Lingkungan: Tingkat deteksi model YOLOv8n dipengaruhi oleh jarak target dan intensitas cahaya. Tingkat deteksi stabil pada 100\% hingga jarak 15 meter, namun turun menjadi 45,5\% pada jarak 20 meter seiring dengan menyusutnya ukuran Bounding Box Area. Selain itu, performa deteksi menurun di bawah 30\% pada kondisi subuh dan sore hari akibat keterbatasan sensor RGB kamera. Pada aspek komunikasi, transmisi aliran video UDP menuju dasbor React rentan mengalami patah-patah (\textit{stuttering}) hingga \textit{freeze} ketika garis pandang terhalang rintangan fisik, meskipun pemrosesan \textit{edge processing} lokal menggunakan Flask dan MAVSDK di dalam drone tetap berjalan secara independen.
\end{enumerate}

Sebagai bagian dari batasan ruang lingkup, penelitian ini belum mengimplementasikan estimasi kebutuhan energi penerbangan yang bersifat adaptif. Sistem belum mempertimbangkan jarak tempuh spasial aktual drone terhadap titik asal (\textit{home position}), belum melakukan prediksi sisa waktu terbang (\textit{remaining flight time}), serta belum mengestimasi konsumsi daya motor propulsi. Fitur keselamatan Battery RTL masih menggunakan nilai ambang batas tetap (\textit{fixed threshold}) sebesar 15\% baterai. Oleh karena itu, pengembangan mekanisme \textit{energy-aware Return to Launch} yang lebih dinamis dapat menjadi arah penelitian lanjutan untuk mengoptimalkan efisiensi energi penerbangan drone otonom.

Secara keseluruhan, penelitian ini mengembangkan sistem deteksi manusia secara otonom yang terintegrasi pada platform \textit{edge computing}. Integrasi model deteksi YOLOv8n yang dioptimalkan dan logika navigasi \textit{Persistent Multi-Target Scouting} terbukti dapat mendukung efisiensi operasional misi pencarian korban. Di samping itu, implementasi Three-Layer Safety System memberikan jaminan redundansi keselamatan terbang guna melindungi drone dari kegagalan operasional. Melalui pencapaian ini, platform terpadu tersebut diharapkan dapat mendukung pengembangan teknologi robotika udara otonom untuk misi kemানুsiaan.

\section{Saran}
\label{sec:sec5_saran}
\vspace{1ex}

Meskipun sistem navigasi otonom dan deteksi manusia pada drone ini telah berfungsi, masih terdapat ruang penyempurnaan untuk pengembangan di masa depan. Berbagai batasan operasional yang teridentifikasi selama pengujian lapangan perlu ditindaklanjuti secara sistematis guna meningkatkan keandalan sistem. Beberapa rekomendasi ilmiah dan teknis disajikan sebagai berikut.

\subsection{Pengembangan Sistem}
\begin{enumerate}
	\item Optimalisasi Manajemen Daya dan Bobot Drone: Mengingat durasi terbang maksimum berkisar antara 4--6 menit, disarankan untuk menambah kapasitas sel baterai (mAh) atau menggunakan konfigurasi baterai paralel. Untuk mengompensasi beban ekstra tersebut, material cetak 3D komponen penyangga drone dapat beralih dari plastik PLA ke material filamen komposit serat karbon (\textit{carbon fiber}) yang lebih ringan namun kokoh guna memperpanjang waktu terbang drone.
	\item Penyematan Akselerator AI dan Modul Penglihatan Malam: Penambahan akselerator perangkat keras (seperti Google Coral Edge TPU atau Hailo-8) disarankan untuk melepaskan beban CPU Raspberry Pi 5 dari inferensi matriks yang berat. Selain itu, pemasangan sensor termal atau inframerah disarankan untuk membuka kapabilitas operasi pencarian pada kondisi minim cahaya.
	\item Mekanisme Kamera \textit{Tilt} dan Transmisi Video Adaptif: Penggunaan kamera monokular tetap (\textit{fixed}) membatasi sapuan sudut pandang. Disarankan menggantinya dengan unit kamera yang memiliki engsel poros \textit{tilt} atau ditopang oleh \textit{brushless gimbal} agar cakupan area observasi lebih dinamis. Secara perangkat lunak, sistem perlu menerapkan algoritma kompresi \textit{Adaptive Bitrate Streaming} untuk meredam hambatan (\textit{stuttering}) video saat terjadi fluktuasi sinyal seluler.
\end{enumerate}

\subsection{Pengembangan Penelitian}
\begin{enumerate}
	\item Implementasi Algoritma \textit{Person Re-Identification} (ReID): Teknik pencegahan penargetan duplikat yang saat ini menggunakan komputasi pengecualian jarak GPS (\textit{Haversine}) sebaiknya disubstitusi dengan visi komputer berupa \textit{Person Re-identification} (ReID). Namun, karena Raspberry Pi 5 memiliki keterbatasan daya komputasi CPU untuk ekstraksi fitur biometrik ganda secara simultan, pengembangan ini disarankan menggunakan perangkat keras \textit{edge computing} dengan akselerator NPU.
	\item Estimasi Spasial dan \textit{Geotagging} Presisi: Penelitian selanjutnya dapat mengembangkan algoritma fotogrametri dinamis untuk memproyeksikan koordinat geografis absolut (Lintang/Bujur) secara otonom dari posisi piksel objek di layar kamera, memadukannya dengan vektor GPS drone, altitudo aktual, dan orientasi \textit{heading yaw}.
	\item Fusi Sensor Ketinggian AGL: Untuk memitigasi fluktuasi pembacaan ketinggian barometer Pixhawk pada operasi ketinggian rendah, penelitian ke depan disarankan menerapkan fusi data (\textit{sensor fusion}) yang menggabungkan pembacaan barometer internal dengan sensor jarak laser LiDAR untuk mengukur ketinggian \textit{Above Ground Level} (AGL) secara presisi.
	\item Pengujian Jaringan Komunikasi Darurat: Diperlukan kajian akademis untuk membandingkan metrik latensi transmisi antara topologi Wi-Fi lokal, satelit mini, dan menara seluler 4G/5G pada berbagai kondisi medan guna merumuskan protokol komunikasi darurat SAR yang tangguh.
\end{enumerate}
